% TOÁN 12 – CHƯƠNG V: PHƯƠNG TRÌNH MẶT PHẲNG, ĐƯỜNG THẲNG, MẶT CẦU TRONG KHÔNG GIAN OXYZ
% PHIẾU BÀI TẬP 01: PHƯƠNG TRÌNH MẶT PHẲNG – MỨC NHẬN BIẾT
\setcounter{subsection}{0}\def\socaulc{12}\def\socauds{2}\def\socaukq{4}\def\socautl{2}

% I. PHÂN DẠNG VÀ VÍ DỤ MINH HỌA
\begin{vd}
Trong không gian với hệ tọa độ $Oxyz$, cho mặt phẳng $(P)\colon x - 2y + 2z - 5 = 0$ và điểm $M(5; 1; c)$. Tìm giá trị của tham số $c$ để điểm $M$ nằm trên mặt phẳng $(P)$.
\loigiai{Vì $M \in (P)$ nên $5 - 2\cdot 1 + 2c - 5 = 0 \Leftrightarrow c = 1$.}
\end{vd}

% II. BÀI TẬP RÈN LUYỆN – MỨC ĐỘ NHẬN BIẾT
\caulc
\Opensolutionfile{ans}[ans/ans-PTMP_Muc1-P1]

\begin{ex}%[2H5N1-3]
Trong không gian $Oxyz$, phương trình mặt phẳng đi qua ba điểm $A(a;0;0)$, $B(0;b;0)$, $C(0;0;c)$ với $abc \ne 0$ là
\choice
{$\dfrac{x}{a} + \dfrac{y}{b} + \dfrac{z}{c} = 0$}
{$\dfrac{x}{b} + \dfrac{y}{c} + \dfrac{z}{a} = 1$}
{\True $\dfrac{x}{a} + \dfrac{y}{b} + \dfrac{z}{c} = 1$}
{$ax + by + cz = 1$}
\loigiai{Phương trình mặt phẳng theo đoạn chắn là $\dfrac{x}{a} + \dfrac{y}{b} + \dfrac{z}{c} = 1$.}
\end{ex}

\begin{ex}%[2H5N1-2]
Cho hình hộp chữ nhật $ABCD.A'B'C'D'$. Bộ bốn vectơ nào sau đây đều là vectơ pháp tuyến của mặt phẳng $(AA'B'B)$?
\choice
{$\overrightarrow{AD}$, $\overrightarrow{A'D'}$, $\overrightarrow{BD}$, $\overrightarrow{B'C'}$}
{$\overrightarrow{AD}$, $\overrightarrow{A'D'}$, $\overrightarrow{BC}$, $\overrightarrow{BC'}$}
{$\overrightarrow{AC}$, $\overrightarrow{A'D'}$, $\overrightarrow{BC}$, $\overrightarrow{B'C'}$}
{\True $\overrightarrow{AD}$, $\overrightarrow{A'D'}$, $\overrightarrow{BC}$, $\overrightarrow{B'C'}$}
\loigiai{Các đường thẳng $AD$, $A'D'$, $BC$, $B'C'$ đều vuông góc với mặt phẳng $(AA'B'B)$.}
\end{ex}

\begin{ex}%[2H5N1-2]
Trong không gian $Oxyz$, mặt phẳng $(P)\colon 2x - 3y + z - 5 = 0$ có một vectơ pháp tuyến là
\choice
{$\overrightarrow{n}_1 = (2; 3; 1)$}
{$\overrightarrow{n}_2 = (2; -3; -5)$}
{$\overrightarrow{n}_3 = (-2; 3; 1)$}
{\True $\overrightarrow{n}_4 = (2; -3; 1)$}
\loigiai{Mặt phẳng $(P)\colon 2x - 3y + z - 5 = 0$ có một vectơ pháp tuyến là $\overrightarrow{n} = (2; -3; 1)$.}
\end{ex}

\begin{ex}%[2H5N1-1]
Trong không gian $Oxyz$, điểm nào dưới đây thuộc mặt phẳng $(P)\colon x - y + 2z - 3 = 0$?
\choice
{$M(1; 1; 1)$}
{$N(2; 1; 0)$}
{\True $P(1; 0; 1)$}
{$Q(0; 1; 1)$}
\loigiai{Thay tọa độ $P(1; 0; 1)$ vào phương trình $(P)$ ta được $1 - 0 + 2\cdot 1 - 3 = 0$.}
\end{ex}

\begin{ex}%[2H5N1-3]
Trong không gian $Oxyz$, mặt phẳng đi qua điểm $M(1; -2; 3)$ và nhận $\overrightarrow{n} = (2; 1; -1)$ làm vectơ pháp tuyến có phương trình là
\choice
{$2x + y - z - 3 = 0$}
{$x - 2y + 3z + 3 = 0$}
{\True $2x + y - z + 3 = 0$}
{$2x - y - z + 3 = 0$}
\loigiai{Phương trình mặt phẳng là $2(x - 1) + 1(y + 2) - 1(z - 3) = 0 \Leftrightarrow 2x + y - z + 3 = 0$.}
\end{ex}

\begin{ex}%[2H5N1-2]
Trong không gian $Oxyz$, mặt phẳng $(Oxy)$ có phương trình là
\choice
{\True $z = 0$}
{$x = 0$}
{$y = 0$}
{$x + y = 0$}
\loigiai{Mặt phẳng $(Oxy)$ có phương trình là $z = 0$.}
\end{ex}

\begin{ex}%[2H5N1-4]
Trong không gian $Oxyz$, hai mặt phẳng $(P)\colon 2x - y + 3z - 1 = 0$ và $(Q)\colon 4x - 2y + 6z + 5 = 0$ có vị trí tương đối là
\choice
{\True Song song}
{Trùng nhau}
{Cắt nhau nhưng không vuông góc}
{Vuông góc}
\loigiai{Ta có $\dfrac{2}{4} = \dfrac{-1}{-2} = \dfrac{3}{6} \ne \dfrac{-1}{5}$ nên $(P) \parallel (Q)$.}
\end{ex}

\begin{ex}%[2H5N1-5]
Trong không gian $Oxyz$, khoảng cách từ điểm $M(1; 2; -1)$ đến mặt phẳng $(P)\colon x + 2y - 2z + 2 = 0$ bằng
\choice
{\True $3$}
{$2$}
{$1$}
{$9$}
\loigiai{Ta có $\mathrm{d}(M, (P)) = \dfrac{|1 + 2\cdot 2 - 2\cdot (-1) + 2|}{\sqrt{1^2 + 2^2 + (-2)^2}} = \dfrac{9}{3} = 3$.}
\end{ex}

\begin{ex}%[2H5N1-2]
Trong không gian $Oxyz$, cho hai vectơ $\overrightarrow{u} = (1; 0; 0)$ và $\overrightarrow{v} = (0; 1; 0)$ là cặp vectơ chỉ phương của mặt phẳng $(P)$. Một vectơ pháp tuyến của mặt phẳng $(P)$ là
\choice
{$\overrightarrow{n} = (1; 1; 0)$}
{$\overrightarrow{n} = (1; 0; 1)$}
{$\overrightarrow{n} = (0; 1; 1)$}
{\True $\overrightarrow{n} = (0; 0; 1)$}
\loigiai{Ta có $\overrightarrow{n} = \left[\overrightarrow{u}, \overrightarrow{v}\right] = (0; 0; 1)$.}
\end{ex}

\begin{ex}%[2H5N1-3]
Trong không gian $Oxyz$, phương trình mặt phẳng đi qua ba điểm $A(2; 0; 0)$, $B(0; -3; 0)$, $C(0; 0; 4)$ là
\choice
{$\dfrac{x}{2} + \dfrac{y}{3} + \dfrac{z}{4} = 1$}
{$\dfrac{x}{2} + \dfrac{y}{-3} + \dfrac{z}{4} = 0$}
{\True $\dfrac{x}{2} + \dfrac{y}{-3} + \dfrac{z}{4} = 1$}
{$2x - 3y + 4z - 1 = 0$}
\loigiai{Phương trình mặt phẳng theo đoạn chắn là $\dfrac{x}{2} + \dfrac{y}{-3} + \dfrac{z}{4} = 1$.}
\end{ex}

\begin{ex}%[2H5N1-4]
Trong không gian $Oxyz$, điều kiện để hai mặt phẳng $(P)\colon A_1x + B_1y + C_1z + D_1 = 0$ và $(Q)\colon A_2x + B_2y + C_2z + D_2 = 0$ vuông góc với nhau là
\choice
{$\dfrac{A_1}{A_2} = \dfrac{B_1}{B_2} = \dfrac{C_1}{C_2}$}
{$A_1A_2 + B_1B_2 + C_1C_2 + D_1D_2 = 0$}
{$A_1B_2 - A_2B_1 = 0$}
{\True $A_1A_2 + B_1B_2 + C_1C_2 = 0$}
\loigiai{Hai mặt phẳng vuông góc khi và chỉ khi $\overrightarrow{n}_1 \cdot \overrightarrow{n}_2 = 0 \Leftrightarrow A_1A_2 + B_1B_2 + C_1C_2 = 0$.}
\end{ex}

\begin{ex}%[2H5N1-5]
Trong không gian $Oxyz$, khoảng cách từ gốc tọa độ $O$ đến mặt phẳng $(P)\colon 2x - 2y + z - 6 = 0$ bằng
\choice
{$6$}
{\True $2$}
{$3$}
{$1$}
\loigiai{Ta có $\mathrm{d}(O, (P)) = \dfrac{|-6|}{\sqrt{2^2 + (-2)^2 + 1^2}} = \dfrac{6}{3} = 2$.}
\end{ex}

\Closesolutionfile{ans}

\cauds
\Opensolutionfile{ans}[ans/ans-PTMP_Muc1-P2]

\begin{ex}%[2H5N1-2]
Trong không gian $Oxyz$, cho mặt phẳng $(P)\colon 2x - y + 3z - 6 = 0$.
\choiceTF
{\True Vectơ $\overrightarrow{n} = (2; -1; 3)$ là một vectơ pháp tuyến của mặt phẳng $(P)$}
{\True Điểm $A(3; 0; 0)$ thuộc mặt phẳng $(P)$}
{Điểm $B(0; 6; 0)$ thuộc mặt phẳng $(P)$}
{Mặt phẳng $(P)$ đi qua gốc tọa độ $O(0; 0; 0)$}
\loigiai{
\begin{itemchoice}
\itemch Mặt phẳng $(P)\colon 2x - y + 3z - 6 = 0$ có một vectơ pháp tuyến là $\overrightarrow{n} = (2; -1; 3)$.
\itemch Thay $A(3; 0; 0)$ vào $(P)$ ta được $2\cdot 3 - 0 + 3\cdot 0 - 6 = 0$.
\itemch Thay $B(0; 6; 0)$ vào $(P)$ ta được $2\cdot 0 - 6 + 3\cdot 0 - 6 = -12 \ne 0$.
\itemch Thay $O(0; 0; 0)$ vào $(P)$ ta được $-6 \ne 0$.
\end{itemchoice}
}
\end{ex}

\begin{ex}%[2H5N1-3]
Xét các mệnh đề viết phương trình mặt phẳng trong không gian $Oxyz$.
\choiceTF
{\True Phương trình mặt phẳng $(P)$ đi qua gốc tọa độ $O(0; 0; 0)$ và có vectơ pháp tuyến $\overrightarrow{n} = (1; -3; 2)$ là $x - 3y + 2z = 0$}
{Phương trình mặt phẳng $(P)$ đi qua $A(1; 2; 3)$, có vectơ pháp tuyến $\overrightarrow{n} = (1; 0; -2)$ là $x - 2z - 5 = 0$}
{\True Phương trình mặt phẳng $(Oxz)$ là $y = 0$ có vectơ pháp tuyến là $(0; 1; 0)$}
{Phương trình mặt phẳng $(P)$ đi qua $M(1; 2; -1)$ và nhận vectơ $\overrightarrow{j}$ làm vectơ pháp tuyến là $z + 1 = 0$}
\loigiai{
\begin{itemchoice}
\itemch Phương trình mặt phẳng đi qua $O(0;0;0)$ có VTPT $\overrightarrow{n}=(1;-3;2)$ là $x - 3y + 2z = 0$.
\itemch Phương trình mặt phẳng là $1(x - 1) + 0(y - 2) - 2(z - 3) = 0 \Leftrightarrow x - 2z + 5 = 0$.
\itemch Mặt phẳng $(Oxz)$ có phương trình $y = 0$, nhận $\overrightarrow{j} = (0; 1; 0)$ làm vectơ pháp tuyến.
\itemch Vì $\overrightarrow{j} = (0; 1; 0)$ nên phương trình mặt phẳng là $1(y - 2) = 0 \Leftrightarrow y - 2 = 0$.
\end{itemchoice}
}
\end{ex}

\Closesolutionfile{ans}

\caukq
\Opensolutionfile{ans}[ans/ans-PTMP_Muc1-P3]

\begin{ex}%[2H5N1-4]
Trong không gian $Oxyz$, cho hai mặt phẳng $(P)\colon 2x - y + 3z + 1 = 0$ và $(Q)\colon -4x + by + cz + 5 = 0$. Khi $(P) \perp (Q)$ thì $3c - b$ bằng bao nhiêu?
\shortans[]{$8$}
\loigiai{Để $(P) \perp (Q)$ thì $2\cdot (-4) + (-1)\cdot b + 3c = 0 \Leftrightarrow 3c - b = 8$.}
\end{ex}

\begin{ex}%[2H5N1-2]
Cho hình hộp chữ nhật $ABCD.A'B'C'D'$. Có bao nhiêu vectơ khác vectơ-không có điểm đầu và điểm cuối là các đỉnh của hình hộp chữ nhật đó và là một vectơ pháp tuyến của mặt phẳng $(ABCD)$?
\shortans[]{$8$}
\loigiai{Gồm $8$ vectơ là $\overrightarrow{AA'}$, $\overrightarrow{A'A}$, $\overrightarrow{BB'}$, $\overrightarrow{B'B}$, $\overrightarrow{CC'}$, $\overrightarrow{C'C}$, $\overrightarrow{DD'}$, $\overrightarrow{D'D}$.}
\end{ex}

\begin{ex}%[2H5N1-2]
Trong không gian $Oxyz$, biết mặt phẳng $(P)\colon x + 6y - 2z + 5 = 0$ có một vectơ pháp tuyến là $\overrightarrow{n} = (a; -6; c)$, trong đó $a$, $c$ là các số nguyên. Tính giá trị của biểu thức $T = ac^2$.
\shortans[]{$-4$}
\loigiai{Mặt phẳng $(P)$ có VTPT $\overrightarrow{n}_0 = (1; 6; -2) \Rightarrow \overrightarrow{n} = -\overrightarrow{n}_0 = (-1; -6; 2)$, suy ra $a = -1$, $c = 2$. Vậy $T = ac^2 = -4$.}
\end{ex}

\begin{ex}%[2H5N1-5]
Trong không gian $Oxyz$, tính khoảng cách từ điểm $M(1; 2; 1)$ đến mặt phẳng $(P)\colon 2x - 2y + z + 4 = 0$.
\shortans[]{$1$}
\loigiai{Ta có $\mathrm{d}(M, (P)) = \dfrac{|2\cdot 1 - 2\cdot 2 + 1 + 4|}{\sqrt{2^2 + (-2)^2 + 1^2}} = 1$.}
\end{ex}

\Closesolutionfile{ans}

\cautl

\begin{ex}%[2H5N1-2]
Cho hình lập phương $ABCD.A'B'C'D'$.
\begin{enumerate}[1.]
\item Tìm một cặp vectơ chỉ phương của mặt phẳng $(ABCD)$.
\item Tìm một vectơ pháp tuyến của mặt phẳng $(ABCD)$.
\end{enumerate}
\loigiai{
\begin{enumerate}[1.]
\item $\overrightarrow{AB}$, $\overrightarrow{AD}$ là một cặp vectơ chỉ phương của mặt phẳng $(ABCD)$.
\item $\overrightarrow{AA'}$ là một vectơ pháp tuyến của mặt phẳng $(ABCD)$.
\end{enumerate}
}
\end{ex}

\begin{ex}%[2H5N1-1]
Trong không gian $Oxyz$, cho hai mặt phẳng $(P)\colon 3x - 5y + 7z + 5 = 0$ và $(Q)\colon x + y - 2 = 0$.
\begin{enumerate}[1.]
\item Tìm một vectơ pháp tuyến của mỗi mặt phẳng $(P)$, $(Q)$.
\item Tìm điểm thuộc mặt phẳng $(P)$ trong số các điểm $A(1; 3; 1)$, $B(1; 2; 3)$.
\end{enumerate}
\loigiai{
\begin{enumerate}[1.]
\item $\overrightarrow{n}_P = (3; -5; 7)$, $\overrightarrow{n}_Q = (1; 1; 0)$.
\item Điểm $A(1; 3; 1)$ thuộc mặt phẳng $(P)$, điểm $B(1; 2; 3)$ không thuộc mặt phẳng $(P)$.
\end{enumerate}
}
\end{ex}
